\section{Discussion}\label{sec:discussion}

\textbf{Which baseline.} A pre-decoder is usually judged against the matching decoder it feeds. Our results show why
that is not enough under realistic noise: most of what the network gains over plain matching at $d = 5$ is recovered,
and more, by correlated matching, a decoding option of the same library. The network's advantage at $d = 3$ and low
error is real but small, and it sits where correlated matching is at its weakest. Any claim of a pre-decoder gain
under device noise should therefore name the strongest network free baseline at the same point. Belief matching
\cite{higgott2023belief} and other correlated decoders were not run and might move the bar further.

\textbf{What calibration buys.} The calibration file helps twice. As a noise source it exposes the network to
structure it was not trained on, and on strongly uneven placements the network's uniform prior costs logical
fidelity. As a set of matching weights it improves plain matching by up to \MedWeightsMax{} times with no network at
all, a gain of the same size as the network's own at $d = 3$. The obvious next step, training or fine tuning the
network on calibrated noise, is outside this study; the released noise model produces exactly the training circuits
it would need.

\textbf{What the card shows.} The card reproduces the integer model exactly, shot for shot in logical outcome, so the
FPGA adds no error beyond quantization. Quantization itself is not free under calibrated noise: the integer pipeline
fails \QuantMin{} to \QuantMax{} times as often as the float one, significantly better at \QuantBetter{} and worse at
\QuantWorse{} of \QuantN{} points. The integer scales are fitted to each circuit's own windows; a deployed integer pre-decoder would need the same per
device calibration. As
an accelerator, one card is slower than the \CpuThreads{} thread host for this small network, and at hundreds of
microseconds per round it is far from a microsecond syndrome cycle. The value of the card in this study is as a verified integer
implementation, not as a speedup.

\section{Limitations}\label{sec:limitations}

No circuit was run on IBM hardware. The calibrations set the rates of a simulated circuit; leakage, crosstalk,
measurement induced dephasing beyond $T_2$, drift and correlated errors are not modelled, and gate errors from
randomized benchmarking are treated as depolarizing. The code is not embedded in the heavy hexagonal lattice, so
the results describe noise drawn from each device's distribution, not that device running the code. The snapshots
are single points in time. NVIDIA's network was used as published, trained at $d = 9$ on uniform noise. Throughput
and latency are measured through PCIe with runtime overheads, not in a decoding loop.
